\documentclass[sigconf,nonacm]{acmart}

%% Pandoc helpers
\providecommand{\tightlist}{%
  \setlength{\itemsep}{0pt}\setlength{\parskip}{0pt}}
\usepackage{booktabs}
\usepackage{longtable}
\usepackage{array}
\usepackage{multirow}
\usepackage{calc}
\usepackage{etoolbox}
\usepackage{textcomp}
\usepackage{amsmath,amssymb}
\usepackage{microtype}
\usepackage{hyperref}

%% ACM bookkeeping
\copyrightyear{2026}
\acmYear{2026}
\setcopyright{rightsretained}
\acmConference[Conference'26]{ACM Conference}{2026}{Venue, City, Country}
\acmBooktitle{Proceedings of ACM Conference (Conference'26)}
\acmISBN{978-x-xxxx-xxxx-x/26/06}
\acmDOI{10.1145/nnnnnnn.nnnnnnn}

\begin{document}

\title{Picking a Workplace Stress Detector by Ablation: An Honest
Model \(\times\) Protocol \(\times\) Sensor \(\times\) Calibration
Evaluation on EmoWork}

\author{Mao Zeth Abel}
\affiliation{%
  \institution{Empathic Project}
  \country{}}
\email{}

\author{Rona Margaret Dalistan}
\affiliation{%
  \institution{Empathic Project}
  \country{}}
\email{}

\author{Jannbeau Amadeus Rain Astrero}
\affiliation{%
  \institution{Empathic Project}
  \country{}}
\email{}

\renewcommand{\shortauthors}{Abel, Dalistan, and Astrero}

\begin{abstract}
We ask a single engineering question: \emph{given a 31-subject
multimodal physiological corpus, which model--protocol--sensor--%
calibration combination is the honest minimum for a deployable
workplace stress detector, and what is the achievable ceiling once a
per-employee calibration step is allowed?} We answer it on
\textbf{EmoWork}, a 31-subject, 12-channel corpus combining cardiac
(ECG, BVP, HR), electrodermal (EDA), thermal (TEMP), inertial (ACC)
and EEG (TP9, AF7, AF8, TP10) recordings during simulated call-centre
knowledge work, with binary stress, binary arousal, binary valence
and four-class affect-quadrant labels. The published multimodal
affect literature routinely reports macro-F1 above 0.80 and
Cohen's \(\kappa\) above 0.60; we show on EmoWork that a substantial
fraction of those numbers are protocol artefacts, that a much smaller
fraction of channels is actually doing the predictive work, and that
almost the entire residual gap to the literature can be closed by
per-subject calibration rather than by architectural change.

The study is structured as a four-axis ablation:
\textbf{model \(\times\) evaluation protocol \(\times\) sensor subset
\(\times\) calibration regime}. We benchmark eleven models---three
classical (Random Forest, Logistic Regression, XGBoost), six deep
fusion learners (1D-CNN, BiLSTM, TinyTCN, Conformer, multi-stream,
DANN-Conformer), one self-supervised pre-trained encoder (TS-TCC),
plus a soft-vote ensemble---under leave-one-subject-out
cross-validation (LOSO). LOSO Cohen's \(\kappa\) peaks at \(0.38\)
for stress (Random Forest), \(0.15\) for arousal (Conformer fusion),
\(0.10\) for valence (BiLSTM fusion) and \(0.05\) for the four-class
quadrant target (XGBoost). To explain the gap, we re-run the three
classical models under three additional protocols: subject-grouped
5-fold, window-stratified 10-fold and a window-stratified 80/20
hold-out. Window-mixing protocols inflate stress \(\kappa\) from 0.31
(LOSO) to \textbf{0.66} (window 80/20); arousal \(\kappa\) jumps from
0.01 to 0.43. Subject-grouped 5-fold tracks LOSO closely, isolating
the inflation to windows of the same subject crossing the train/test
boundary. A per-modality ablation reframes the multimodal narrative:
ECG alone (17 features) equals the full 149-feature fusion stack on
stress, and adding EEG to an 89-d physio stack reduces arousal
\(\kappa\) from 0.08 to 0.02. A within-subject 70/30 calibration
regime (\(\approx\!14\) labelled windows per subject) lifts macro-F1
to 0.91 (stress, LogReg), 0.82 (arousal, BiLSTM), 0.86 (valence,
BiLSTM) and 0.72 (quadrant, CNN1D), exceeding DEAP/AMIGOS
within-subject baselines by \(\approx\!+0.25\). The contribution is a
defensible four-axis floor and ceiling: macro-F1 \(\approx 0.70\) for
the generalisation problem, \(\approx 0.91\) for the personalisation
problem, with explicit accounting for how every \(+0.05\) above
either threshold in the literature is most plausibly explained by
protocol leakage or by silent per-subject evaluation. Per-subject
calibration, not model architecture, is the dominant lever.
\end{abstract}

\begin{CCSXML}
<ccs2012>
  <concept>
    <concept_id>10010147.10010178.10010179</concept_id>
    <concept_desc>Computing methodologies~Machine learning</concept_desc>
    <concept_significance>500</concept_significance>
  </concept>
  <concept>
    <concept_id>10010520.10010553.10010562</concept_id>
    <concept_desc>Computer systems organization~Embedded and cyber-physical systems</concept_desc>
    <concept_significance>300</concept_significance>
  </concept>
  <concept>
    <concept_id>10003120.10003123.10010860</concept_id>
    <concept_desc>Human-centered computing~User studies</concept_desc>
    <concept_significance>300</concept_significance>
  </concept>
</ccs2012>
\end{CCSXML}

\ccsdesc[500]{Computing methodologies~Machine learning}
\ccsdesc[300]{Computer systems organization~Embedded and cyber-physical systems}
\ccsdesc[300]{Human-centered computing~User studies}

\keywords{EmoWork, multimodal affect recognition, leave-one-subject-out,
data leakage, modality ablation, per-subject calibration,
personalisation, electrocardiography, electrodermal activity,
electroencephalography, stress detection, valence, arousal,
Cohen's kappa}

\maketitle

$body$

\end{document}
